% Generated draft: requires mathematical and source review.
\documentclass[11pt]{article}
\usepackage[margin=25mm]{geometry}
\usepackage{fontspec}
\setmainfont{Latin Modern Roman}
\usepackage{amsmath,amssymb,mathtools,unicode-math}
\setmathfont{Latin Modern Math}
\usepackage{enumitem,needspace,xurl,hyperref}
\setlength{\parindent}{0pt}
\setlength{\parskip}{5pt}
\setlength{\emergencystretch}{3em}
\newcommand{\N}{\mathbb{N}}
\newcommand{\Z}{\mathbb{Z}}
\newcommand{\Q}{\mathbb{Q}}
\newcommand{\R}{\mathbb{R}}
\newcommand{\C}{\mathbb{C}}
\newcommand{\F}{\mathbb{F}}
\newcommand{\A}{\mathbb{A}}
\newcommand{\PP}{\mathbb{P}}
\newcommand{\E}{\mathbb{E}}
\newcommand{\eps}{\varepsilon}
\newcommand{\dd}{\,\mathrm{d}}
\begin{document}
\section*{UnsolvedMath ID: 30005349}
\textbf{Source problem code:} catalogue entry
\subsection{Definitions and mathematical statement}
An indicator function on a finite set \(\Omega\) is a function \(f:\Omega\to\{0,1\}\). Under uniform averaging, for every nonempty subset \(T\subseteq\Omega\), define
\[
\mathbb E_T f=\frac{1}{|T|}\sum_{x\in T}f(x).
\]
The supplied record does not explicitly specify the averaging measure. Under the customary Boolean-coordinate interpretation, the balanced slice would be
\[
\Omega_n=\left\{x\in\{0,1\}^n:\sum_{i=1}^n x_i=n/2\right\},\qquad n\in2\N.
\]
This ambient-domain convention requires confirmation; it is not defined in the supplied text.

Let \(\mathcal S_n\) denote the intended family of nonempty sub-balanced slices of the ambient balanced slice, and let \(d(T)\) denote their intended dimension. These are placeholders: the supplied record does not define either the permitted sub-balanced slices or their dimension.

\textbf{Question, conditional on resolving these definitions.} For every indicator function \(f\) on the balanced slice satisfying \(\mathbb E_{\Omega_n}f=1/2\), does there exist \(T\in\mathcal S_n\) such that
\[
d(T)=\log^*n
\qquad\text{and}\qquad
\mathbb E_Tf>0.51?
\]
Here \(\log^*n\) is retained from the source. The usual base-two convention is the number of successive base-two logarithms needed to obtain a value at most \(1\); the record does not specify a convention, rounding adjustment, or admissible dimension restriction.

The source also asks for the weaker conclusion that the dimension tends to infinity. A precise possible asymptotic reading of that alternative is: there exist an integer-valued function \(g(n)\to\infty\) and a threshold \(N\), both independent of \(f\), such that for every admissible \(n\ge N\) and every such \(f\), there exists \(T\in\mathcal S_n\) with
\[
d(T)\ge g(n),\qquad \mathbb E_Tf>0.51.
\]
This quantifier interpretation, the admissible values of \(n\), and the threshold convention require confirmation. The record does not determine a fully specified mathematical statement without these qualifications.
\subsection{Short English statement}
The supplied question asks whether every Boolean function of mean \(1/2\) on a balanced slice has a sub-balanced slice on which its mean exceeds \(0.51\), with dimension \(\log^*n\), or at least with dimension growing without bound. The permitted sub-slices, dimension convention, averaging measure, and precise asymptotic quantifiers need clarification.
\subsection{Sources}
\begin{itemize}
\item Local UnsolvedMath registry, supplied complete record: permanent ID \(30005349\), title \emph{Biased Sub-Balanced Slices of Boolean Functions}, exact source problem code \texttt{catalogue entry}. The record belongs to the Oberwolfach Reports Open Problems collection and lists proposed year \(2023\); no proposer is supplied.
\item Supplied report citation: \emph{Combinatorics (2023)}, Oberwolfach Reports, DOI \texttt{10.4171/owr/2023/1}, pp.\ 5--89. Source URL and supplied literature-source link: \url{https://doi.org/10.4171/owr/2023/1}.
\item Supplied literature assessment, checked exactly \texttt{2026-08-22}: the proposed sub-balanced-slice strengthening is described as open in the source report, and the supplied curation reports finding no primary resolution by that date. This is supplied dated curation, not an independently verified assessment. The background explicitly requests independent verification of time-sensitive status and cited claims.
\item The registry labels the statement \texttt{corrected\_verified} and says it was checked against the official report. That verification claim and the supplied \texttt{open} status are recorded as source metadata; neither has been independently verified here.
\end{itemize}
\end{document}
